% SPDX-License-Identifier: Apache-2.0
% covers: I2c
\datasheet{I2c}{A general I2C master on AXI-Lite: one piece of a transaction per command, on two open drain lines.}

\dsfacts{\code{lib/parts/src/i2c.rs}}{\code{txhdl\_parts::i2c::I2c}}%
{\code{//docs:examples}}

\subsection*{Function}

\code{I2cInit} of the HDMI part replays a fixed list of writes and
cannot read. This master is the other kind: a program says what one
piece of a transaction is and the master does it, so a driver can
address a device, write to it, turn the bus around and read from it.

One command is at most a start, one byte either way, and a stop. A
driver that wants to read a register writes the address with
\code{start|write}, then asks for \code{start|read|nack}, whose start
is the repeated start. Between the commands of one transaction the
lines are parked with the clock low, since a data line that moves
while the clock is high is a start or a stop to every device on the
bus.

\subsection*{Register map}

Offsets from the block's base.

\dsregs{I2c}

A write to \code{cmd} starts the command it describes, and \code{cmd}
reads zero. A write of ones to \code{state} clears \code{fired},
\code{nack} and \code{lost}; \code{busy} is read only.

A quarter of a bit takes \code{div + 1} cycles, so the bus runs at the
system clock over \code{4 * (div + 1)}. From 100 MHz, 249 is
100 kbit/s and 62 is 400 kbit/s.

Bit 4 of \code{cmd} is set for a NACK, which is what a master sends
after the last byte it wants; a master that acknowledges the last byte
tells the device to send another.

\dstables{I2c}

\subsection*{The two lines}

Both are open drain, as I2C requires: a device pulls a line low and
nobody drives one high. \code{scl\_low} and \code{sda\_low} high pull
their line low, and a line nobody pulls is high through the board's
pull-up, so a board wrapper drives a pad from each output and reads
the pad back into \code{scl\_in} and \code{sda\_in}.

The master watches both inputs, and two things follow that a master
driving the lines itself could not do.

A device that holds the clock low is stretching it, and the master
waits for the line to rise rather than counting cycles. And a data
line that reads low while this master released it is another master
winning the bus, which sets arbitration lost and ends the command.

\subsection*{Behaviour}

The master is two processes. The bus process answers the registers
every cycle, counts the cycles of a quarter of a bit in \code{tick},
drives \code{scl\_low} and \code{sda\_low} from \code{scl\_pull} and
\code{sda\_pull}, and latches what a command found. The engine is the
transaction written as the sequence it is: a wait for a command, then
the start, the eight bits, the acknowledge and the stop, a quarter of
a bit per wait, with the lines set for each quarter before its wait.
The lowering numbers the waits into a state register the unit does not
declare, \code{at\_wait}, and the eight bits are a counted loop with a
counter of its own, \code{for0}; both are in the tables above, and
there is no \code{step} or \code{quarter} register to read.

\begin{itemize}
\item Writing \code{cmd} while no command runs records what was asked
and the byte to send, clears \code{nack\_hit} and \code{lost\_hit},
sets \code{busy}, and sets \code{held} on a start.
\item A quarter ends when \code{tick} reaches the divider and the
clock is not being stretched: a device holding the clock low where
the master released it holds the quarter with it.
\item On a write, the acknowledge reads the data line: low is the
device acknowledging, high sets \code{nack\_hit}.
\item On a read, the byte arrives in \code{shift} and is put in
\code{data}, which the word at \code{0x08} reads. The master then
drives the answer bit 4 asked for.
\item A master writing a one that reads back a zero has lost the bus
to another master. It sets \code{lost\_hit}, releases both lines for
the rest of the byte, whose bits the loop still counts through, and
skips the acknowledge and the stop.
\item At the end the lines are parked, the clock low while the
transaction stays open and both released when it is over or the bus
was lost. \code{done} is up for one cycle, when the bus process
latches \code{fired}, \code{nack} and \code{lost}; \code{busy} falls a
cycle after it. \code{fired} is the done bit, which raises the
interrupt while the enable is set.
\item Out of reset both lines are released: \code{scl\_pull} and
\code{sda\_pull} say whether the engine pulls a line, and reset to
not pulling (issue 606).
\end{itemize}

\subsection*{Verification}

\begin{itemize}
\item \code{//lib/parts:parts\_test} drives the master against the
device model in the same file: a write reaches the device, a read
takes a byte from it, an address nobody answers sets not acknowledged,
a device that stretches the clock is waited for, and another master
that wins the third bit of the address makes this one report
arbitration lost and let go of the data line.
\item \code{//lib/examples:ex\_i2c} addresses a device, writes to it,
turns the bus around and reads back, and the build simulates the
netlist against that run under nvc and Verilator.
\end{itemize}

\subsection*{Used by}

Nothing in the tree yet. The HDMI part still uses \code{I2cInit} for
its fixed setup; issue 148 notes that \code{I2cInit} could become a
program for this master instead, and that a board's EEPROM, which
often holds the Ethernet address, is what wants a master that can
read.

\subsection*{Limits}

Ten-bit addressing is not here: an address is a byte like any other,
so a driver writes what it likes, but nothing helps it. There is no
general call handling, no clock synchronisation with another master
beyond losing arbitration, and no timeout on a stretched clock, so a
device that holds the clock low for ever holds the master with it.
